\section{对偶间隙何时闭合}
\label{sec:09-Convex-Optimization/05-Duality/02}

你为一个最小化问题构造了对偶函数，跑了一轮数值求解。原始目标值是 \(p=10.4\)，对偶目标值是 \(d=10.1\)。两个数字之间的 0.3 意味着什么？它可能只是数值误差，也可能是一个真实的、理论上不可消除的鸿沟。

弱对偶只保证方向：

\[
d^\star\le p^\star.
\]

它没有说这两个值之间距离多大。对偶间隙正式定义就是这个差：

\[
p^\star-d^\star.
\]

若间隙为零，即

\[
d^\star=p^\star,
\]

称为强对偶。强对偶不只是``两个问题碰巧答案一致''。它为原始问题提供了一份全局下界证书：不仅找到了一个解，还证明了不可能有更好的解。这意味着如果某天有人拿着一个更低的原始目标值过来，你可以直接告诉他：要么他的点不可行，要么他算错了——因为对偶值已经认证了不可能更低。

\subsection{凸性加可行不等于间隙为零}

线性规划的强对偶非常稳：在适当的可行和有界条件下，间隙自动为零。把``LP 总有强对偶''这个经验带到一般凸问题上，以为``凸 + 可行''就是强对偶的充分条件——会有反例打脸。

在定义域 \(y>0\) 上考虑：

\[
\begin{aligned}
\text{minimize}\quad&e^{-x}\\
\text{subject to}\quad&\frac{x^2}{y}\le0.
\end{aligned}
\]

函数 \(x^2/y\) 是二次除线性函数，在 \(y>0\) 上的 Hessian 为

\[
\nabla^2\left(\frac{x^2}{y}\right)=\frac{2}{y^3}\begin{bmatrix}y^2 & -xy \\ -xy & x^2\end{bmatrix}.
\]

对于任意向量 \([u,v]^{\trans}\)，二次型 \(u^2 y^2-2uvxy+v^2 x^2/(y^3)=(uy-vx)^2/(y^3)\ge0\)，因此 Hessian 半正定——函数凸。约束 \(x^2/y\le0\) 在 \(y>0\) 时强制 \(x^2\le0\)，即必须 \(x=0\)。可行域只有一条半直线 \(\{(0,y):y>0\}\)。代入目标：\(p^\star=e^0=1\)。

构造 Lagrangian：

\[
L(x,y,\lambda)=e^{-x}+\lambda\frac{x^2}{y},\qquad\lambda\ge0.
\]

对任意固定的 \(\lambda\ge0\)，对偶函数为

\[
g(\lambda)=\inf_{x,\,y>0}L(x,y,\lambda).
\]

关键的技巧在于 \(y\)：固定 \(x\)，让 \(y\to\infty\)，惩罚项 \(\lambda x^2/y\to0\)。因此

\[
g(\lambda)=\inf_x e^{-x}=\lim_{x\to\infty}e^{-x}=0.
\]

对所有 \(\lambda\ge0\) 取上确界：\(d^\star=\sup_{\lambda\ge0} g(\lambda)=0\)。结果：

\[
d^\star=0<1=p^\star.
\]

问题凸。可行域非空。目标下方有界。间隙却严格为正。

问题的根在于：约束 \(x^2/y\le0\) 在可行域上永远等于零，不存在任何让约束严格小于零的点。可行域没有``内部''——它是一条线，紧贴在约束边界上。几何分离需要把超平面放在可行域的``内部一侧''，但在这里根本没有内部可放。这正是约束资格条件要排除的退化情形。

\subsection{Slater 条件：让可行域有一口呼吸的空间}

对标准形式凸问题

\[
\begin{aligned}
\min_x\quad&f_0(x)\\
\text{s.t.}\quad&f_i(x)\le0,\\
&Ax=b,
\end{aligned}
\]

Slater 条件要求：存在一个点 \(\tilde x\) 在共同定义域的相对内部，满足

\[
f_i(\tilde x)<0
\quad\text{对所有非仿射不等式},
\qquad
A\tilde x=b.
\]

入门教材通常给一个更强的版本：存在满足全部等式且让每一个不等式都严格小于零的点。对纯仿射不等式约束（不需要严格成立）可在相对内部上做更细致的处理；对于第一次遇到这个概念的读者，记住关键画面：必须在不等式围成的区域内部找到一个点，而不是只能在边界上找。

用二维画面建立几何直觉。想象约束围成一个圆盘。如果可行域包含圆心（严格可行点），对偶函数在最优乘子处的水平集就是一条支撑整个圆盘的线——强对偶成立。如果可行域只有圆周（没有内部点），那就相当于一个圆形``线框''——分离超平面仍然存在，但法向量的方向退化了。Slater 条件就是在说：那个圆不能只是一个圆圈，里面得有点东西。

为什么内部点的存在就能让对偶间隙闭合？对偶理论的标准证明构造了一个集合

\[
\mathcal{A}=\{(u,v,t)\mid\exists x,\;f_0(x)\le t,\;f_i(x)\le u_i,\;Ax-b=v\},
\]

即所有可达到的（约束值，目标值）组合。如果 \(p^\star\) 以下的值不可达，则 \((0,0,p^\star)\) 不在 \(\mathcal{A}\) 的内部。由凸集的分离超平面定理，存在一个非零法向量分离该点与 \(\mathcal{A}\)。这个法向量的分量就是 Lagrange 乘子 \((\lambda,\nu)\)。

分离定理只能保证存在一个非零法向量——却不一定能保证 \(\lambda\) 的各个分量可以被归一化为合法的、非负的乘子。Slater 条件的作用，就是确保 \(\mathcal{A}\) 在约束值的方向上具有足够的``厚度''（至少能取到一组严格负的约束值），使得法向量指向\(\lambda>0\) 的那个方向。

通俗地说：凸性保证了集合的外形没有凹陷（分离定理的前提成立），Slater 条件保证了分离出的超平面是可用的证书（法向量分量非负、非全零）。

\subsection{Slater 是充分条件而非必要条件}

一个问题不满足 Slater，不表示强对偶一定失败。很多边界问题间隙仍然为零，只是不能由 Slater 这条定理提供保证。标准的线性规划就是一个例子：它的约束全部是仿射不等式和等式，Slater 条件可以放宽为``存在可行点''即可（严格内部不是必须的）。另一方面，原问题非凸也可能偶然出现零间隙——但这是巧合，无法系统依赖。

正确的推理路径是：

\begin{itemize}
\tightlist
\item
  若问题凸且满足 Slater：强对偶成立，且对偶最优可取得（在适当闭性条件下）。
\item
  若问题凸但 Slater 失败：需要检查其他约束资格——如只有线性不等式约束时可放松条件、广义 Slater 条件（对非仿射不等式仍要求严格内部点）、或直接分析该问题的几何结构。
\item
  ``没找到严格可行点''不等于``不存在严格可行点''——没找到可能是搜索不充分或建模时就把边界写死了。
\item
  一次数值实验中原始值与对偶值很接近，也不能反推所有理论条件已满足——数值接近可能来自算法早停、问题缩放不当或浮点截断。
\end{itemize}

\subsection{最优值相等不等于最优点存在}

\(p^\star=d^\star\) 只是在说两个实数值相等。它没有保证：

\begin{itemize}
\tightlist
\item
  存在原始可行点 \(x^\star\) 使得 \(f_0(x^\star)=p^\star\)（原始最优取得到）；
\item
  存在对偶可行乘子 \((\lambda^\star,\nu^\star)\) 使得 \(g(\lambda^\star,\nu^\star)=d^\star\)（对偶最优取得到）。
\end{itemize}

两者可能最优值有限，却只能由一列点逼近：例如 \(\inf_{x>0} 1/x = 0\)，但不存在 \(x>0\) 使得 \(1/x=0\)。\(x_k\to\infty\) 时目标趋近下确界但永远达不到。实分析中 ``infimum attained'' 和 ``infimum not attained'' 的区分在这里具有实际后果：你的求解器可能永远无法返回一个真正最优的解，只会在真值附近来回振荡——因为那个真值本身就是达不到的。

实际中需要分开确认五项：

\begin{itemize}
\tightlist
\item
  原始可行点是否存在；
\item
  原始最优值是否被某个可行点取得；
\item
  对偶可行乘子是否存在；
\item
  对偶最优值是否被某个可行乘子取得；
\item
  间隙是否为零。
\end{itemize}

在满足 Slater 的凸问题中，前三项可由定理保证（对偶最优也可取得，但需声明闭性条件），但每一项都应查阅具体的定理陈述，而不是打包成一个模糊的``满足条件就都好''。

\subsection{对偶间隙是停止条件的最强证据}

回到开头的场景：你手上有原始可行点 \(\hat x\)（目标值 10.4），对偶可行乘子 \((\hat\lambda,\hat\nu)\)（对偶值 10.1）。无论强对偶是否被理论证明，弱对偶已经给出

\[
10.1\le p^\star\le10.4.
\]

这个区间不是由启发式规则给出的——不是``梯度已经小于 \(10^{-5}\) 所以差不多收敛了''，不是``过去 100 步 loss 只降了 0.1\%''。它是一个严格的数学界限：真实最优值一定在这个范围内。\(\hat x\) 的绝对次优度不超过 0.3——你最多还可以把目标压低 0.3，不能再多。

这个区间的宽度本身就是停止证据。很多优化算法用``迭代变化量很小''作为停止信号。那个信号告诉你的是``算法走不动了''，但对偶间隙告诉你的是``离最优还有多远''——前者可能是走到了一个平坦区而并非真正最优（比如被 saddle point 骗住了），后者是一个硬数学边界。两者的可信度不在一个量级上。

若还满足强对偶条件，上下界最终收敛到同一个数，则两侧同时获得认证：不仅知道了 \(p^\star\) 是多少，还拥有了一份下界证明，解释了为什么不可能存在比它更好的解。

对于目标值尺度差异大的问题（一个问题的最优值是 \(10^{-3}\)，另一个是 \(10^6\)），绝对间隙没有可比性，使用相对间隙：

\[
\frac{p^\star-d^\star}{|p^\star|+\varepsilon}.
\]

分母加 \(\varepsilon\) 防止 \(p^\star\approx 0\) 时的数值问题。注意不同求解器报告的 \texttt{gap} 定义并不统一——有的是绝对间隙，有的是相对间隙，有的对分母做了其他归一化。读文档比信任一个无注释的浮点数更可靠。

对偶间隙最终要回答的问题很简单：我现在的解离最优还有多远，以及我能为``不可能更好了''提供什么级别的数学保证。弱对偶给出的是单向不等式，强对偶把它升级为等号。从不等式到等号的跨越，正是 Slater 条件在几何上确保的那一口呼吸。
